% =========================================================
% Cardinality and Infinite Sets
% =========================================================

\subsection{Cardinality and Infinite Sets}

\begin{exposition}
The intuitive idea of size works perfectly for finite sets: finite sets can be
compared by counting. For infinite sets, counting is replaced by the existence
of functions with the right structural properties. Cantor's central idea is
that two sets have the same size exactly when their elements can be matched by
a bijection.
\end{exposition}

% ---------------------------------------------------------
% Same Cardinality
% ---------------------------------------------------------

\begin{definitionbox}{Definition (Same Cardinality)}
\begin{definition}[Same Cardinality]\label{def:same-cardinality}
Two sets $A$ and $B$ have the \emph{same cardinality}, written $A \sim B$,
if there exists a bijection $f : A \to B$.
\end{definition}
\end{definitionbox}

\begin{remark*}[Standard quantified statement]
$A\sim B \leftrightarrow \exists f:A\to B,\;
\operatorname{Bijective}(f,A,B)$.
\end{remark*}

\begin{remark*}[Predicate reading]
$\operatorname{SameCardinality}(A,B)$ means that there is a bijection from
$A$ onto $B$.
\end{remark*}

\begin{remark*}[Interpretation]
Same cardinality means perfect matchability with no repetitions and no missed
elements.
\end{remark*}

\begin{remark*}[Exposition]
The relation $\sim$ is an equivalence relation on sets. Reflexivity is
witnessed by the identity function, symmetry by inverses of bijections, and
transitivity by composition of bijections. The equivalence collection of $A$ under
$\sim$ is called the cardinality of $A$, written $|A|$.
\end{remark*}

\begin{remark*}[Exposition]
Let $E=\{2,4,6,\dots\}$ be the set of positive even integers. The function
$f:\mathbb{N}\to E$ defined by $f(n)=2n$ is a bijection, so
$\mathbb{N}\sim E$. Thus an infinite set can have the same cardinality as a
proper subset of itself.
\end{remark*}

\begin{remark*}[Exposition]
Dedekind's characterization says that an infinite set can be matched
bijectively with some proper subset of itself. This phenomenon cannot occur
for finite sets.
\end{remark*}

\begin{dependencies}
\begin{itemize}
  \item \hyperref[def:bijective]{Bijective Function}
  \item \hyperref[def:identity]{Identity Function}
  \item \hyperref[def:composition]{Composition}
\end{itemize}
\end{dependencies}

% ---------------------------------------------------------
% Finite, Countable, Uncountable
% ---------------------------------------------------------

\begin{definitionbox}{Definition (Finite and Infinite Sets)}
\begin{definition}[Finite and Infinite Sets]\label{def:finite-infinite}
A set $A$ is \emph{finite} if $A=\varnothing$ or
$A\sim\{1,2,\dots,n\}$ for some $n\in\mathbb{N}$. A set is \emph{infinite}
if it is not finite.
\end{definition}
\end{definitionbox}

\begin{remark*}[Standard quantified statement]
$\operatorname{FiniteSet}(A)\leftrightarrow
(A=\varnothing\vee \exists n\in\mathbb{N},\;
\operatorname{SameCardinality}(A,\{1,\dots,n\}))$, and
$\operatorname{InfiniteSet}(A)\leftrightarrow\neg\operatorname{FiniteSet}(A)$.
\end{remark*}

\begin{remark*}[Predicate reading]
$\operatorname{FiniteSet}(A)$ means that $A$ is empty or bijective with an
initial segment of $\mathbb{N}$; $\operatorname{InfiniteSet}(A)$ means that
this finite condition fails.
\end{remark*}

\begin{remark*}[Interpretation]
Finite sets are exactly the sets whose elements can be exhausted by a finite
count.
\end{remark*}

\begin{dependencies}
\begin{itemize}
  \item \hyperref[def:empty-set]{Empty Set}
  \item \hyperref[def:same-cardinality]{Same Cardinality}
  \item \hyperref[def:natural-numbers]{Natural Numbers}
\end{itemize}
\end{dependencies}

\begin{definitionbox}{Definition (Countable and Uncountable Sets)}
\begin{definition}[Countable and Uncountable Sets]\label{def:countable}
A set $A$ is \emph{countably infinite} if $A\sim\mathbb{N}$. A set is
\emph{countable} if it is finite or countably infinite. A set is
\emph{uncountable} if it is infinite but not countable.
\end{definition}
\end{definitionbox}

\begin{remark*}[Standard quantified statement]
$\operatorname{CountablyInfiniteSet}(A)\leftrightarrow
\operatorname{SameCardinality}(A,\mathbb{N})$;
$\operatorname{CountableSet}(A)\leftrightarrow
\operatorname{FiniteSet}(A)\vee \operatorname{CountablyInfiniteSet}(A)$; and
$\operatorname{UncountableSet}(A)\leftrightarrow
\operatorname{InfiniteSet}(A)\wedge \neg\operatorname{CountableSet}(A)$.
\end{remark*}

\begin{remark*}[Predicate reading]
$\operatorname{CountableSet}(A)$ means that $A$ is finite or can be listed by
natural-number indices.
\end{remark*}

\begin{remark*}[Interpretation]
Countability means listability, allowing either a finite list or an infinite
list indexed by $\mathbb{N}$.
\end{remark*}

\begin{remark*}[Exposition]
A countably infinite set admits a listing $a_1,a_2,a_3,\dots$ with every
element appearing exactly once. Such a listing is a bijection
$\mathbb{N}\to A$.
\end{remark*}

\begin{dependencies}
\begin{itemize}
  \item \hyperref[def:finite-infinite]{Finite and Infinite Sets}
  \item \hyperref[def:same-cardinality]{Same Cardinality}
  \item \hyperref[def:natural-numbers]{Natural Numbers}
\end{itemize}
\end{dependencies}

% ---------------------------------------------------------
% Ordering cardinalities
% ---------------------------------------------------------

\begin{definitionbox}{Definition (Comparing Cardinalities)}
\begin{definition}[Comparing Cardinalities]\label{def:cardinality-order}
Write $|A|\leq |B|$ if there exists an injection $f:A\hookrightarrow B$.
Write $|A|<|B|$ if $|A|\leq |B|$ but $A\not\sim B$.
\end{definition}
\end{definitionbox}

\begin{remark*}[Standard quantified statement]
$\operatorname{CardinalityLeq}(A,B)\leftrightarrow
\exists f:A\to B,\;\operatorname{Injective}(f,A,B)$, and
$\operatorname{CardinalityLt}(A,B)\leftrightarrow
\operatorname{CardinalityLeq}(A,B)\wedge
\neg\operatorname{SameCardinality}(A,B)$.
\end{remark*}

\begin{remark*}[Predicate reading]
$\operatorname{CardinalityLeq}(A,B)$ means that an injective copy of $A$ fits
inside $B$.
\end{remark*}

\begin{remark*}[Interpretation]
Cardinal comparison is expressed by embeddings, not by attempting to count
infinite sets directly.
\end{remark*}

\begin{remark*}[Exposition]
The direction matters. An injection $A\hookrightarrow B$ places a copy of
$A$ inside $B$, while an injection $B\hookrightarrow A$ gives the opposite
comparison.
\end{remark*}

\begin{dependencies}
\begin{itemize}
  \item \hyperref[def:injective]{Injective Function}
  \item \hyperref[def:same-cardinality]{Same Cardinality}
\end{itemize}
\end{dependencies}

% ---------------------------------------------------------
% Countability theorems
% ---------------------------------------------------------

\begin{theorembox}{Theorem ($\mathbb{Q}$ is countable)}
\begin{theorem}[$\mathbb{Q}$ is countable]\label{thm:Q-countable}
The set $\mathbb{Q}$ of rational numbers is countable.
\hyperref[prf:Q-countable]{\textit{Go to proof.}}
\end{theorem}
\end{theorembox}

\begin{remark*}[Standard quantified statement]
$\operatorname{CountableSet}(\mathbb{Q})$.
\end{remark*}

\begin{remark*}[Predicate reading]
$\operatorname{CountableSet}(\mathbb{Q})$ means that the rational numbers can
be listed by natural-number indices.
\end{remark*}

\begin{remark*}[Interpretation]
Although $\mathbb{Q}$ is dense in $\mathbb{R}$, it is still listable.
\end{remark*}

\begin{remark*}[Exposition]
The proof arranges positive rationals as an infinite two-dimensional array and
reads the array along diagonals, skipping duplicates. Negative rationals and
zero are then interleaved into the same list.
\end{remark*}

\begin{dependencies}
\begin{itemize}
  \item \hyperref[def:countable]{Countable Set}
  \item \hyperref[def:natural-numbers]{Natural Numbers}
  \item \hyperref[def:bijective]{Bijective Function}
\end{itemize}
\end{dependencies}

\begin{theorembox}{Theorem (Countable Union of Countable Sets)}
\begin{theorem}[Countable Union of Countable Sets]\label{thm:countable-union}
Let $\{A_n\}_{n\in\mathbb{N}}$ be a countable family of countable sets.
Then $\bigcup_{n=1}^{\infty}A_n$ is countable.
\hyperref[prf:countable-union]{\textit{Go to proof.}}
\end{theorem}
\end{theorembox}

\begin{remark*}[Standard quantified statement]
$\forall (A_n)_{n\in\mathbb{N}},\;
(\forall n\in\mathbb{N},\operatorname{CountableSet}(A_n))\to
\operatorname{CountableSet}(\bigcup_{n=1}^{\infty}A_n)$.
\end{remark*}

\begin{remark*}[Predicate reading]
If each $A_n$ is countable, then the union indexed by $\mathbb{N}$ is also
countable.
\end{remark*}

\begin{remark*}[Interpretation]
Countability is stable under countable unions.
\end{remark*}

\begin{remark*}[Exposition]
The proof uses a diagonal enumeration. Arrange listings of the sets $A_n$ as
rows of an infinite array and read the array along anti-diagonals, ignoring
repetitions.
\end{remark*}

\begin{dependencies}
\begin{itemize}
  \item \hyperref[def:countable]{Countable Set}
  \item \hyperref[def:natural-numbers]{Natural Numbers}
  \item \hyperref[def:indexed-family]{Indexed Family of Sets}
  \item \hyperref[def:union]{Union}
\end{itemize}
\end{dependencies}

% ---------------------------------------------------------
% Uncountability
% ---------------------------------------------------------

\begin{theorembox}{Theorem ($\mathbb{R}$ is uncountable)}
\begin{theorem}[$\mathbb{R}$ is uncountable]\label{thm:R-uncountable}
The set $\mathbb{R}$ of real numbers is not countable.
\hyperref[prf:R-uncountable]{\textit{Go to proof.}}
\end{theorem}
\end{theorembox}

\begin{remark*}[Standard quantified statement]
$\operatorname{UncountableSet}(\mathbb{R})$, equivalently
$\neg\operatorname{CountableSet}(\mathbb{R})$.
\end{remark*}

\begin{remark*}[Predicate reading]
$\operatorname{UncountableSet}(\mathbb{R})$ means that no natural-number-indexed
list contains every real number.
\end{remark*}

\begin{remark*}[Interpretation]
The real numbers have strictly larger cardinality than any countable set.
\end{remark*}

\begin{remark*}[Exposition]
Cantor's diagonal argument defeats any proposed list
$r_1,r_2,r_3,\dots$ by constructing a real number whose $n$-th digit differs
from the $n$-th digit of $r_n$.
\end{remark*}

\begin{remark*}[Exposition]
Since $\mathbb{Q}$ is countable and $\mathbb{R}$ is uncountable, the
irrational numbers are uncountable. Otherwise $\mathbb{R}$ would be a
countable union of countable sets.
\end{remark*}

\begin{dependencies}
\begin{itemize}
  \item \hyperref[def:countable]{Countable Set}
  \item \hyperref[def:finite-infinite]{Finite and Infinite Sets}
  \item \hyperref[thm:countable-union]{Countable Union of Countable Sets}
\end{itemize}
\end{dependencies}

% ---------------------------------------------------------
% Cantor's theorem
% ---------------------------------------------------------

\begin{theorembox}{Theorem (Cantor's Theorem)}
\begin{theorem}[Cantor's Theorem]\label{thm:cantor}
For any set $A$, there is no surjection $A\to\mathcal{P}(A)$. In particular,
$|A|<|\mathcal{P}(A)|$.
\hyperref[prf:cantor]{\textit{Go to proof.}}
\end{theorem}
\end{theorembox}

\begin{remark*}[Standard quantified statement]
$\forall A,\;\forall f:A\to\mathcal{P}(A),\;
\neg\operatorname{Surjective}(f,A,\mathcal{P}(A))$.
\end{remark*}

\begin{remark*}[Predicate reading]
No function from $A$ to $\mathcal{P}(A)$ hits every subset of $A$.
\end{remark*}

\begin{remark*}[Interpretation]
The power set of any set is strictly larger than the original set.
\end{remark*}

\begin{remark*}[Exposition]
Given $f:A\to\mathcal{P}(A)$, form
\[
D=\{a\in A\mid a\notin f(a)\}.
\]
The set $D$ differs from every candidate value $f(a)$ at the element $a$, so
$D$ is not in the image of $f$.
\end{remark*}

\begin{remark*}[Exposition]
Cantor's theorem implies there is no largest cardinality. Starting from any
set $A$, the power set $\mathcal{P}(A)$ has strictly larger cardinality.
\end{remark*}

\begin{dependencies}
\begin{itemize}
  \item \hyperref[def:function]{Function}
  \item \hyperref[def:surjective]{Surjective Function}
  \item \hyperref[def:power-set]{Power Set}
  \item \hyperref[ax:separation]{Axiom Schema of Separation}
\end{itemize}
\end{dependencies}

% ---------------------------------------------------------
% Characteristic functions and B^A
% ---------------------------------------------------------

\begin{definitionbox}{Definition (Function Space $B^A$)}
\begin{definition}[Function Space $B^A$]\label{def:function-space}
For sets $A$ and $B$, the \emph{function space} $B^A$ denotes the set of all
functions from $A$ to $B$:
\[
B^A:=\{f:A\to B\}.
\]
\end{definition}
\end{definitionbox}

\begin{remark*}[Standard quantified statement]
$g\in B^A\leftrightarrow g:A\to B$.
\end{remark*}

\begin{remark*}[Predicate reading]
$\operatorname{FunctionSpace}(B^A,A,B)$ means that $B^A$ is the set of all
functions with domain $A$ and codomain $B$.
\end{remark*}

\begin{remark*}[Interpretation]
Function spaces package all functions of a fixed type into a single set.
\end{remark*}

\begin{remark*}[Exposition]
For finite sets with $|A|=m$ and $|B|=n$, the number of functions $A\to B$ is
$n^m$, since there are $n$ choices for each of the $m$ inputs.
\end{remark*}

\begin{remark*}[Exposition]
For any set $A$, the power set $\mathcal{P}(A)$ and the function space
\(\{0,1\}^A\) have the same cardinality via the characteristic function map
$S\mapsto \mathbf{1}_S$, where $\mathbf{1}_S(a)=1$ if $a\in S$ and
$\mathbf{1}_S(a)=0$ if $a\notin S$.
\end{remark*}

\begin{dependencies}
\begin{itemize}
  \item \hyperref[def:function]{Function}
  \item \hyperref[def:set-membership]{Set and Membership}
  \item \hyperref[def:power-set]{Power Set}
\end{itemize}
\end{dependencies}

% ---------------------------------------------------------
% Schr\"{o}der--Bernstein
% ---------------------------------------------------------

\begin{theorembox}{Theorem (Schr\"{o}der--Bernstein Theorem)}
\begin{theorem}[Schr\"{o}der--Bernstein Theorem]\label{thm:schroder-bernstein}
If $|A|\leq |B|$ and $|B|\leq |A|$, then $A\sim B$.
Equivalently, if there exist injections $f:A\hookrightarrow B$ and
$g:B\hookrightarrow A$, then there exists a bijection $h:A\to B$.
\hyperref[prf:schroder-bernstein]{\textit{Go to proof.}}
\end{theorem}
\end{theorembox}

\begin{remark*}[Standard quantified statement]
$\operatorname{CardinalityLeq}(A,B)\wedge\operatorname{CardinalityLeq}(B,A)
\to \operatorname{SameCardinality}(A,B)$.
\end{remark*}

\begin{remark*}[Predicate reading]
Two opposite cardinal injections force a bijection between the two sets.
\end{remark*}

\begin{remark*}[Interpretation]
To prove two sets have the same cardinality, it is enough to embed each one
into the other.
\end{remark*}

\begin{remark*}[Exposition]
Schr\"{o}der--Bernstein is useful when a direct bijection is hard to describe.
The theorem turns two easier injection constructions into an existence theorem
for a bijection.
\end{remark*}

\begin{remark*}[Exposition]
A sequence in a set $X$ is a function $\mathbb{N}\to X$. Thus sequence
notation is a special case of function-space notation.
\end{remark*}

\begin{remark*}[Exposition]
The distinction between countable and uncountable sets reappears throughout
analysis, including separability, countable unions in measure theory, and
function spaces.
\end{remark*}

\begin{dependencies}
\begin{itemize}
  \item \hyperref[def:cardinality-order]{Comparing Cardinalities}
  \item \hyperref[def:injective]{Injective Function}
  \item \hyperref[def:bijective]{Bijective Function}
\end{itemize}
\end{dependencies}
